\documentclass[12pt]{article}
\usepackage[utf8]{inputenc}
\usepackage[russian]{babel}
\usepackage{amsmath}
\usepackage{amssymb}
\usepackage{geometry}
\geometry{a4paper, margin=2cm}

\begin{document}

\section*{Задача 3.1.4}

\textbf{Условие:} Вывести двоичное представление десятичного числа. Least significant bit справа (0.5 балла).\\
in: x (1 регистр) \\
out: y (log2(x) регистров)

\section*{Решение}

\begin{verbatim}
Read 1  ; x

Load 1
JZero END_0     ; x == 0 => Out: 0

WHILE: Load =0
    Sub 1
    JGTZ END
    JZero END   ; exit if x <= 0
    Load 1
    Div =2
    Mult =2
    Sub 1
    Mult =-1    ; r0 := x mod 2 = ((x // 2) * 2 - x) * -1
    Write 0
    Load 1
    Div =2
    Store 1     ; x := x // 2
    Jump WHILE
END_0: Write =0
END: Halt 0
\end{verbatim}

\section*{Сложность}

Тут мы не рассматриваем $x < 0$ (на паре сказали, что не надо и что это можно понять по тому, что нужно вывести log2(x) регистров). Худший случай - $x > 0$, тогда цикл будет выполняться столько раз, сколько бит в двоичном представлении x. Пусть \(n := \lfloor \log_2 x \rfloor + 1\).

\begin{enumerate}
    \item Инициализация за 3 шага
    \item Цикл WHILE выполняется n раз
    \item Каждая успешная итерация выполняется за 14 шагов
    \item После n-й итерации выполняется ещё одна проверка условия - 4 шага
    \item Halt 0 - 1 шаг
\end{enumerate}

Значит сложность в худшем случае:
\[
    T_{\max} = 3 + 14n + 4 + 1 = 14n + 8
\]
где $n = \lfloor \log_2 x \rfloor + 1$

\end{document}